% !TeX program = lualatex
% !TeX encoding = utf8
% !TeX spellcheck = uk_UA
% !TeX root = ../main.tex



%=========================================================
\Opensolutionfile{answer}[\currfilebase/\currfilebase-Answers]
\Writetofile{answer}{\protect\section*{\nameref*{\currfilebase}}}
\chapter{Енергія магнітного поля}\label{\currfilebase}\hypertarget{MagEnergy}{}
\graphicspath{{\currfilebase/Pictures}}
%=========================================================
\epigraph{\Alexander Енергія поля існує у~двох різних формах, які можна описати без гіпотез як магнітну поляризацію та
електричну поляризацію або, за дуже ймовірною гіпотезою, як рух і~напруження одного й~того самого середовища.}{Джеймс Клерк
Максвелл, \enquote{Динамічна теорія електромагнітного поля} (1865)}

На відміну від енергії електростатичного поля, яку можна ввести відразу після закону
Кулона (просто обчисливши роботу перенесення зарядів проти кулонівської сили), з~магнітною
енергією так просто не вийде: магнітна складова сили Лоренца, з~якою поле діє на заряд,
роботи не виконує (розділ~\ref{MagFieldVacuum})~--- за жодних обставин. Виникає природне
запитання: якщо магнітне поле роботи не виконує, звідки ж тоді береться енергія, яку воно,
як виявиться, запасає?

Відповідь дає ЕРС \emphx{самоіндукції}~--- окремий випадок закону електромагнітної
індукції (розділ~\ref{ElMagInduction}): джерело, що підтримує струм у~контурі, витрачає
частину роботи не на подолання магнітної сили (вона роботи не виконує), а~на подолання
самої цієї ЕРС, яку контур наводить сам у~собі через власне змінюване поле. Тому
послідовно ввести енергію магнітного поля можна лише \emph{після} того, як встановлено
закон індукції і~поняття самоіндукції; сам механізм~--- чому магнітна заборона тут не
діє і~хто насправді виконує цю роботу~--- описано в~\S~\ref{sec:Work_current}, а~підсумок підведено наприкінці розділу.

%% --------------------------------------------------------
\section{Енергія магнітного поля}\hypertarget{magnetic-field-energy}{}
%% --------------------------------------------------------



%% --------------------------------------------------------
\subsection*{Енергія контуру зі струмом}
%% --------------------------------------------------------



Провідник зі струмом створює магнітне поле, яке з'являється й~зникає разом зі струмом;
\emph{магнітне поле~--- носій енергії}, що дорівнює роботі джерела струму на його
створення. Розглянемо роботу джерела ЕРС $\emf$ під час замикання ключа в~колі з~котушкою
індуктивності $L$ й~резистором $R$, при перенесенні заряду $dq = I\, dt$. За законом Ома $\emf+\emf_\text{сі}=IR$, де $\emf_\text{сі}=-\dfrac{L}{c^2}\dfrac{dI}{dt}$~--- ЕРС самоіндукції~\eqref{eq:SelfInductionEMF}, тому
\begin{equation*}
	\delta A = \emf\, dq = I^2R\, dt - \emf_\text{сі}I\, dt = \underbrace{I^2R\, dt}_\text{теплота} +
	\underbrace{\frac1{c^2}LI\, dI}_{\text{енергія поля}}.
\end{equation*}
Частина роботи джерела йде на теплоту (закон Джоуля--Ленца), а~решта~--- на створення
магнітного поля. Інтегруючи другий доданок від нульового струму до кінцевого значення $I$,
отримуємо енергію магнітного поля контуру:
\begin{equation}\label{eq:MagEnergyFormula}
	W = \frac1{c^2}\int\limits_0^I LI\, dI = \frac{LI^2}{2c^2} = \frac{I\Phi}{2c} = \frac{\Phi^2}{2L}.
\end{equation}



%% --------------------------------------------------------
\subsection*{Об'ємна густина енергії магнітного поля: соленоїд}
%% --------------------------------------------------------



Знайдемо\index{Густина енергії!магнітного поля} енергію магнітного поля, використовуючи характеристики самого поля, на прикладі
однорідного поля довгого соленоїда з~проникністю $\mu$
(рис.~\ref{pic:SolenoidCrossSection}). Індуктивність соленоїда~--- за
формулою~\eqref{eq:SolenoidInductance}, $L = 4\pi\mu N^2 S/\ell$, а~поле всередині~--- $B
= 4\pi\mu NI/(c\ell)$ (розділ~\ref{MagFieldMedia}). Підставляючи
в~\eqref{eq:MagEnergyFormula}:

\begin{SCfigure}[3][h!]\centering
	\localinput{solenoid-cross-section.tikz}
	\caption{Поздовжній переріз довгого соленоїда: однорідне поле всередині, струм
	у~витках обмотки\label{pic:SolenoidCrossSection}}
\end{SCfigure}

\begin{equation*}
	W = \frac1{c^2}\frac{LI^2}2 = \frac1{c^2}\cdot \frac{(4\pi\mu NI)^2 S\ell}{8\pi\mu \ell^2} = \frac{B^2}{8\pi\mu}\, V = \frac{BH}{8\pi}\, V,
\end{equation*}
де $V=S\ell$~--- об'єм соленоїда (об'єм, зайнятий полем), а~в~останній рівності
використано $H=B/\mu$. Звідси об'ємна густина енергії магнітного поля:
\begin{equation}\label{eq:MagEnergyDensitySolenoid}
	\tcbhighmath{
		w = \frac{BH}{8\pi}.
	}
\end{equation}

Як і~для електричного поля (розділ~\ref{ElEnergy}), \emph{магнітна енергія зосереджена саме в~тому об'ємі простору, де є магнітне поле}: у~випадку соленоїда це об'єм усередині нього, а~не сам дріт обмотки.

Формула~\eqref{eq:MagEnergyDensitySolenoid} справедлива не лише для соленоїда: у~квазістаціонарному наближенні її можна вивести
із закону індукції для довільного розподілу поля (див. нижче), а~строго, з~урахуванням струму зміщення,~--- з~теореми
Пойнтінга (\S~\ref{sec:PintingTheoreme}). Для лінійного середовища ($\Bfield=\mu\Hfield$) вона є точною магнітною паралеллю до густини енергії електричного поля $ED/8\pi$ (розділ~\ref{ElEnergy}):
\begin{equation}\label{eq:MagEnergyDensityGeneral}
	\tcbhighmath{
		W = \iiint\limits_V \frac{\Bfield\cdot\Hfield}{8\pi}\, dV, \qquad
		w = \frac{\Bfield\cdot\Hfield}{8\pi} = \frac{\mu H^2}{8\pi}.
	}
\end{equation}


%% --------------------------------------------------------
\subsection*{Густина енергії в~загальному випадку*}
%% --------------------------------------------------------
\begin{Extra}
	Покажемо, що формула~\eqref{eq:MagEnergyDensitySolenoid} не обмежується соленоїдом. Обмежимося \emph{квазістаціонарним}
	наближенням: поля змінюються повільно (розділ~\ref{SteadyCurrent}), тож струмом зміщення нехтуємо, $\divg\vect j=0$
	і~$\rot\Hfield=\frac{4\pi}{c}\vect j$.

	Магнітне поле роботи над зарядами не виконує (розділ~\ref{MagFieldVacuum}); роботу над носіями струму виконує вихрове електричне
	поле $\Efield$, яке виникає, коли магнітне поле змінюється. Потужність, яку поле $\Efield$ передає струмам
	в~одиниці об'єму, дорівнює $\vect j\cdot\Efield$, тому за час $dt$
	\begin{equation*}
		\delta A = dt\iiint\limits_V \vect j\cdot\Efield\, dV = \frac{c\, dt}{4\pi}\iiint\limits_V \Efield\cdot\rot\Hfield\, dV.
	\end{equation*}
	За тотожністю $\divg(\Efield\times\Hfield) = \Hfield\cdot\rot\Efield - \Efield\cdot\rot\Hfield$ і~теоремою
	Остроградського--Гаусса
	\begin{equation*}
		\delta A = \frac{c\,dt}{4\pi}\iiint\limits_V \Hfield\cdot\rot\Efield\, dV
		- \frac{c\,dt}{4\pi}\oiint\limits_{S} (\Efield\times\Hfield)\cdot d\vect S.
	\end{equation*}
	Для обмеженої системи у~квазістаціонарному наближенні поверхневий інтеграл по нескінченно віддаленій поверхні дорівнює нулю
	(випромінювання немає). У~першому доданку $\rot\Efield = -\frac1c\,\partial\Bfield/\partial t$~\eqref{eq:CurlEDiff}:
	\begin{equation*}
		\delta A = -\frac{dt}{4\pi}\iiint\limits_V \Hfield\cdot\frac{\partial\Bfield}{\partial t}\, dV.
	\end{equation*}
	Це робота вихрового поля над струмами; вона від'ємна: поле \emph{забирає} енергію у~струмів (точніше, у~джерел, що їх
	підтримують) і~передає її магнітному полю. Тому приріст магнітної енергії $dW=-\delta A$:
	\begin{equation}\label{eq:MagEnergyDifferential}
		dW = \frac1{4\pi}\iiint\limits_V \Hfield\cdot d\Bfield\, dV.
	\end{equation}
	Цей вираз справедливий і~для нелінійних середовищ: для феромагнетика за цикл перемагнічування він дає втрати
	$\frac1{4\pi}\oint H\,dB$ в~одиниці об'єму, тобто площу петлі гістерезису (розділ~\ref{MagFieldMedia}). Для лінійного
	середовища ($\Bfield=\mu\Hfield$) $\Hfield\cdot d\Bfield = \tfrac12 d(\Hfield\cdot\Bfield)$, і~інтегрування
	дає~\eqref{eq:MagEnergyDensityGeneral}.

	Строго, з~урахуванням струму зміщення, цей результат разом з~енергією електричного поля дає теорема Пойнтінга
	(розділ~\ref{MaxwellEqns}, \S~\ref{sec:PintingTheoreme}).
\end{Extra}



%% --------------------------------------------------------
\section{Взаємна індукція}\hypertarget{mutual-inductance}{}
%% --------------------------------------------------------



%% --------------------------------------------------------
\subsection*{Коефіцієнти взаємної індуктивності}
%% --------------------------------------------------------



Нехай два нерухомі контури 1 і~2 розташовані близько один до одного
(рис.~\ref{pic:TwoContours}). Якщо в~контурі 1 тече струм $I_1$, він створює крізь контур
2 магнітний потік, пропорційний $I_1$:
\begin{equation}\label{eq:MutualFlux21}
	\Phi_{21} = \frac1c L_{21}I_1;
\end{equation}
аналогічно, струм $I_2$ у~контурі 2 створює крізь контур 1 потік
\begin{equation}\label{eq:MutualFlux12}
	\Phi_{12} = \frac1c L_{12}I_2.
\end{equation}

\begin{SCfigure}[1.25][h!]\centering
	\localinput{12contours.tikz}
	\caption{Два нерухомі контури, магнітно пов'язані один з~одним\label{pic:TwoContours}}
\end{SCfigure}

Коефіцієнти пропорційності $L_{12}$ і~$L_{21}$ називають \emphx{коефіцієнтами взаємної
індуктивності} контурів; на відміну від власної індуктивності $L$, яка завжди
додатна (енергія $LI^2/2c^2>0$), це \emph{алгебраїчні} величини, які можуть мати будь-який знак (залежно від
взаємної орієнтації контурів), а~також дорівнювати нулю (якщо контури взаємно
перпендикулярні так, що жоден потік іншого через себе не пронизує). Знак $L_{12}$ залежить і~від вибору додатних напрямків обходу: для двох співвісних кілець, обійдених в~одному напрямку, потоки додаються й~$L_{12}>0$, а~за протилежного обходу одного з~них $L_{12}$ змінює знак.

Повний магнітний потік, що пронизує кожен із контурів, складається з~власного потоку
(самоіндукції) та потоку, створеного струмом іншого контуру:
\begin{equation}\label{eq:MutualFluxes}
	\begin{cases}
		\Phi_1 = \Phi_{11} + \Phi_{12}
		= \dfrac1c\,L_{11}I_1 + \dfrac1c\,L_{12}I_2, \\[1.5ex]
		\Phi_2 = \Phi_{21} + \Phi_{22}
		= \dfrac1c\,L_{21}I_1 + \dfrac1c\,L_{22}I_2,
	\end{cases}
\end{equation}
де $\Phi_{11}=\dfrac1c L_{11}I_1$ і~$\Phi_{22}=\dfrac1c L_{22}I_2$ --- власні потоки,
а~$\Phi_{12}$ і~$\Phi_{21}$ задані співвідношеннями \eqref{eq:MutualFlux12}
і~\eqref{eq:MutualFlux21} відповідно.

З~урахуванням взаємної індукції закони Ома для двох контурів набувають вигляду
\begin{equation}\label{eq:MutualOhm}
	\begin{cases}
		I_1R_1 = \emf_1 - \dfrac1{c^2}\dfrac{d}{dt}\left(L_{11}I_1\right)
		- \dfrac1{c^2}\dfrac{d}{dt}\left(L_{12}I_2\right), \\[1.5ex]
		I_2R_2 = \emf_2 - \dfrac1{c^2}\dfrac{d}{dt}\left(L_{21}I_1\right)
		- \dfrac1{c^2}\dfrac{d}{dt}\left(L_{22}I_2\right),
	\end{cases}
\end{equation}
де $\emf_1,\emf_2$~--- ЕРС джерел у~контурах, $L_{11}\equiv L_1$, $L_{22}\equiv L_2$~--- власні індуктивності контурів. Той факт, що
зміна струму в~одному контурі неминуче породжує ЕРС і~в~іншому, і~називають \emphx{явищем
взаємної індукції}.

У~загальному випадку для \(N\) нерухомих контурів повний магнітний потік, що пронизує
\(i\)-й контур, складається з~власного потоку та потоків, створених струмами всіх інших
контурів, і~записується через суму:
\begin{equation}\label{eq:FluxSumN}
	\Phi_i = \sum_{j=1}^{N} \Phi_{ij}
	= \frac1c \sum_{j=1}^{N} L_{ij} I_j,
	\qquad i = 1, 2, \dots, N,
\end{equation}
де діагональні коефіцієнти \(L_{ii}\)~--- власні індуктивності контурів, а~недіагональні
\(L_{ij}\) (\(i\neq j\))~--- \emphx{коефіцієнти взаємної індуктивності}\index{Індуктивність!взаємна} \(i\)-го та
\(j\)-го контурів.

Підставляючи \eqref{eq:FluxSumN} у~закон електромагнітної індукції (ЕРС індукції в~\(i\)-му контурі дорівнює \(-\dfrac1c\dfrac{d\Phi_i}{dt}\)), отримуємо закон Ома для цього контуру, де \(\emf_i\)~--- ЕРС джерела в~ньому:
\begin{equation}\label{eq:OhmN}
	I_i R_i = \emf_i - \frac1{c^2}\frac{d}{dt}\sum_{j=1}^{N} L_{ij} I_j,
	\qquad i = 1, 2, \dots, N.
\end{equation}

\begin{Example}[unbreakable]{Дві обмотки на спільному осерді}
	Довгий парамагнітний циліндр довжиною $\ell$ і~площею перерізу $S$ (проникність $\mu$) несе дві обмотки, намотані одна поверх
	одної по всій довжині: перша
	з~густиною намотування $n_1$ витків на одиницю довжини, друга --- $n_2$; знайдемо їхню
	взаємну індуктивність, нехтуючи крайовими ефектами (рис.~\ref{pic:TwoCoils}).

	Оскільки $L_{21}=c\Phi_2/I_1$, потрібно створити струм $I_1$ в~обмотці 1 і~обчислити
	повний потік крізь усі $N_2 = n_2\ell$ витків обмотки 2 (аналогічно $N_1=n_1\ell$): $\Phi_2 = N_2 B_1 S$, де $B_1
	= \dfrac{4\pi}{c}\mu n_1 I_1$ (розділ~\ref{MagFieldMedia}), тому
	\begin{equation}\label{eq:MutualExample}
		L_{21} = 4\pi\mu\, \frac{N_1N_2}{\ell}S,
	\end{equation}
	і, обчисливши аналогічно $L_{12}$, отримуємо той самий результат:
	$L_{12}=L_{21}=4\pi\mu N_1N_2S/\ell$.
\end{Example}

\begin{figure}[h!]\centering
	\includegraphics[width=0.6\linewidth]{two_coils}
	\caption{Дві обмотки на спільному осерді: обмотка~1 створює поле, що пронизує витки обмотки~2\label{pic:TwoCoils}}
\end{figure}

\begin{Attention}[unbreakable]
	У~цьому прикладі $L_{12}=L_{21}=\sqrt{L_1L_2}$. Ця гарна рівність~--- наслідок двох
	спрощень: (1)~ми знехтували розсіюванням магнітного потоку поза осердям,
	і~(2)~проникність осердя $\mu$ не залежить від величини поля. Для феромагнітного
	осердя, де $\mu=\mu(H)$, і~за умови $N_1\neq N_2$, обмотки при однакових струмах
	створюють у~осерді \emph{різні} поля, тож вирази для $L_{21}$ і~$L_{12}$ міститимуть
	різні (польозалежні) значення $\mu$, і~рівність $L_{21}=L_{12}$, узагалі кажучи,
	порушиться. Диференціальні ж коефіцієнти $\partial\Phi_2/\partial I_1$ і~$\partial\Phi_1/\partial I_2$ збігаються завжди (за відсутності гістерезису), бо енергія є функцією стану.
\end{Attention}



%% --------------------------------------------------------
\subsection*{Теорема взаємності}
%% --------------------------------------------------------



У~щойно розглянутому прикладі коефіцієнти $L_{12}$ і~$L_{21}$ обчислювалися двома
геометрично різними шляхами, а~зійшлися в~одному й~тому самому значенні. Це не
випадковість прикладу з~високосиметричним осердям: насправді рівність
\begin{equation}\label{eq:MutualEqual}
	\tcbhighmath{
		L_{12} = L_{21}.
	}
\end{equation}
виконується для \emph{будь-якої} пари контурів довільної форми
(за єдиної умови~--- лінійності середовища, $\mu\neq f(H)$, як і~вказано вище). Тобто,
результати розрахунку потоку в~обох напрямках точно збігаються. Це твердження називається
\emphx{теоремою взаємності}\index{Теорема взаємності}.

Практично корисний наслідок теореми: потік, який струм $I$ контуру 1 наводить крізь контур 2, дорівнює потоку, який такий самий струм $I$ контуру 2 навів би крізь контур 1, $\Phi_{12}=\Phi_{21}$. Тому немає потреби розрізняти $L_{12}$ і~$L_{21}$~--- говоритимемо про \emphx{взаємну індуктивність двох контурів} $L_{12}$. Відповідно ЕРС, яку змінний струм $I(t)$ контуру 2 наводить у~контурі 1, дорівнює ЕРС, яку той самий струм, якби він ішов контуром 1, навів би в~контурі 2.

Чому теорема має бути правильною, видно з~енергії. Увімкнемо спочатку струм $I_1$ за відсутності $I_2$, а~потім збільшимо $I_2$ від нуля до кінцевого значення за $I_1=\const$. Змінний струм $I_2$ наводить у~контурі~1 ЕРС $-\dfrac{L_{12}}{c^2}\dfrac{dI_2}{dt}$; джерело контуру~1 її долає й~за весь процес виконує додаткову роботу $L_{12}I_1I_2/c^2$. Якщо ввімкнути струми в~іншому порядку, дістанемо $L_{21}I_1I_2/c^2$. Енергія стану з~готовими струмами $I_1$, $I_2$ не залежить від того, як цей стан створено (інакше цикл \enquote{вмикання~--- вимикання} давав би вічний двигун), отже $L_{12}=L_{21}$. Строгі доведення~--- нижче.

Доведемо цю теорему. Нехай маємо два нерухомі тонкі контури $C_1$ та $C_2$. За означенням,
струм $I_2$ у~контурі $C_2$ створює крізь контур $C_1$ магнітний потік
(рис.~\ref{fig:mutual-inductance-reciprocity})
\begin{equation}\label{eq:MutualFlux12_def}
	\Phi_{12} = \frac{1}{c} L_{12} I_2.
\end{equation}

%---------------------------------------------------------
\begin{SCfigure}[1.75][h!]\centering

	\localinput{mutual-inductance-reciprocity.tikz}
	\caption{До доведення теореми взаємності. Зображено контури $C_1$ і~$C_2$, елементи
	$d\vect{l}_1$, $d\vect{l}_2$ та радіус-вектор $r_{21}$ між ними, які використовуються
	при обчисленні потоку $\Phi_{12}$.}
	\label{fig:mutual-inductance-reciprocity}
\end{SCfigure}
%---------------------------------------------------------

Запишемо потік $\Phi_{12}$ крізь поверхню $S_1$, що спирається на контур $C_1$, від поля
$\Bfield^{(2)}$, створюваного струмом $I_2$. Поклавши $\Bfield^{(2)} = \rot\vect A^{(2)}$
(див.~\eqref{eq:BrotA}), де $\vect A^{(2)}$~--- векторний потенціал поля другого струму, за
теоремою Стокса маємо:
\begin{equation}\label{eq:FluxViaVectorPotential}
	\Phi_{12} = \iint\limits_{S_1}\rot\vect A^{(2)}\cdot d\vect S_1 = \oint\limits_{C_1}\vect A^{(2)}\cdot d\vect l_1,
\end{equation}
де $d\vect l_1$~--- орієнтований елемент контуру $C_1$.

Для тонкого провідника $C_2$ векторний потенціал $\vect A^{(2)}$ у~точці розташування
елемента $d\vect l_1$ виражається контурним інтегралом~\eqref{eq:VectorPotentialDef} (для
лінійних струмів $\vect{j} dV \to I d\vect{l}$):
\begin{equation}\label{eq:VectorPotentialLineCurrent}
	\vect A^{(2)} = \frac{\mu I_2}{c}\oint\limits_{C_2}\frac{d\vect l_2}{r_{12}},
\end{equation}
де $r_{12}$~--- відстань між елементами $d\vect l_2$ та $d\vect l_1$.

Підставляючи вираз~\eqref{eq:VectorPotentialLineCurrent}
у~формулу~\eqref{eq:FluxViaVectorPotential}, отримуємо подвійний контурний інтеграл:
\begin{equation}\label{eq:FluxDoubleIntegral}
	\Phi_{12} = \frac{\mu I_2}{c}\oint\limits_{C_1}\oint\limits_{C_2}\frac{d\vect l_1\cdot d\vect l_2}{r_{12}}.
\end{equation}

Порівнюючи вираз~\eqref{eq:FluxDoubleIntegral} із вихідним означенням
потоку~\eqref{eq:MutualFlux12_def}, знаходимо явний вигляд взаємної індуктивності (формулу
Неймана):
\begin{equation}\label{eq:MutualInductanceNeumann}
	L_{12} = {\mu}\oint\limits_{C_1}\oint\limits_{C_2}\frac{d\vect l_1\cdot d\vect l_2}{r_{12}}.
\end{equation}

Проводячи повністю симетричні міркування для коефіцієнта $L_{21}$ (обчислюючи потік
$\Phi_{21} = \frac{1}{c}L_{21}I_1$ від струму $I_1$ крізь контур $C_2$), дістаємо:
\begin{equation}\label{eq:MutualInductanceNeumann21}
	L_{21} = {\mu}\oint\limits_{C_2}\oint\limits_{C_1}\frac{d\vect l_2\cdot d\vect l_1}{r_{21}}.
\end{equation}

Оскільки скалярний добуток елементів $d\vect l_1 \cdot d\vect l_2 = d\vect l_2 \cdot
d\vect l_1$ є комутативним, а~відстань між ними $r_{12} = r_{21}$ однакова в~обох
випадках, підінтегральний вираз у~формулі Неймана є повністю симетричним відносно
перестановки індексів $1 \leftrightarrow 2$. Звідси безпосередньо випливає теорема
взаємності: \(	L_{12} = L_{21} \).

Формула~\eqref{eq:MutualInductanceNeumann} прозоро ілюструє суто \emph{геометричний}
характер взаємної індуктивності: величина $L_{12}$ залежить лише від форми, розмірів та
взаємного розташування двох контурів у~просторі, але абсолютно не залежить від величин
струму в~них. Знак внеску легко простежити: співнапрямлені елементи, $d\vect l_1\cdot d\vect l_2>0$, дають у~$L_{12}$ внесок зі знаком плюс, зустрічні~--- мінус. Це доведення проведено для однорідного середовища ($\mu=\const$ у~всьому просторі); для довільного лінійного середовища, зокрема для обмоток на осерді, теорему доведено нижче через інтеграл по об'єму (підрозділ про вираз енергії через поля).

\begin{figure}[tbp]\centering
	\localinput{coaxial-loops.tikz}
	\caption{Два коаксіальні колові контури різних радіусів\label{pic:CoaxialLoops}}
\end{figure}

\begin{Example}{Два коаксіальні кільця}
	Два тонкі колові провідники зі спільною віссю лежать в~одній площині; радіус
	зовнішнього~--- $R_1$, внутрішнього~--- $R_2$, причому $R_2\ll R_1$
	(рис.~\ref{pic:CoaxialLoops}). Знайдемо магнітний потік крізь контур зовнішнього
	провідника, якщо внутрішнім тече струм $I$.

	Пряме обчислення потоку від маленького кільця крізь велике~--- складна геометрична
	задача (поле маленького кільця істотно неоднорідне на масштабі великого). Натомість
	застосуємо \emphx{теорему взаємності}: обчислимо потік у~\emph{зворотному}
	напрямку~--- від великого кільця крізь маленьке,~--- це набагато простіше, оскільки на
	масштабі маленького кільця поле великого можна вважати практично однорідним і~рівним
	полю в центрі великого кільця, $B_1 = 2\pi I/(cR_1)$ (розділ~\ref{MagFieldVacuum}). Тоді
	\begin{equation*}
		\Phi_{21} = B_1\cdot \pi R_2^2 = \frac{2\pi^2 I R_2^2}{cR_1},
	\end{equation*}
	і, за теоремою взаємності, той самий потік (при тому самому струмі $I$) виникає
	й~у~прямій задачі:
	\begin{equation*}
		\Phi_{12} = \Phi_{21} = \frac{2\pi^2 I R_2^2}{cR_1}.
	\end{equation*}
\end{Example}

\begin{figure}[tbp]\centering
	\localinput{rotating-dipole.tikz}
	\caption{Магнітний диполь, що обертається навколо осі $\omega$, і~нерухоме кільце радіуса $a$ на відстані $\ell$; $\vect n$~--- нормаль до кільця\label{pic:RotatingDipole}}
\end{figure}

\begin{Example}{Обертовий магнітний диполь}
	Магнітний диполь з~моментом $p_m$ обертається з~циклічною частотою $\omega$ навколо осі, що проходить через його центр
	перпендикулярно до $\vect p_m$. На відстані $\ell\gg a$ від диполя, на прямій у~площині його обертання, розташоване нерухоме
	плоске колове кільце радіуса $a$ і~опору $R$; нормаль $\vect n$ кільця напрямлена вздовж цієї прямої
	(рис.~\ref{pic:RotatingDipole}). Самоіндукцією кільця знехтуємо (це можна, якщо $\omega L/c^2\ll R$). Знайдемо струм
	у~кільці.

	Оскільки $\vect n$ напрямлено на диполь, нормальна до кільця складова поля~\eqref{eq:DipoleFieldB}
	(розділ~\ref{MagFieldVacuum}) дорівнює $B_n=2\,\vect p_m\!\cdot\vect n/\ell^3$. Диполь обертається в~площині, що містить
	$\vect n$, тож $\vect p_m\cdot\vect n=p_m\cos\omega t$. Через $\ell\gg a$ поле майже однорідне на площі кільця, і
	\begin{equation*}
		\Phi(t)=\pi a^2B_n=\frac{2\pi a^2p_m}{\ell^3}\cos\omega t.
	\end{equation*}
	Звідси ЕРС індукції $\emf=-\frac1c\,\dfrac{d\Phi}{dt}=\dfrac{2\pi a^2p_m\omega}{c\,\ell^3}\sin\omega t$, а~струм
	$I(t)=\emf/R$ гармонічно змінюється з~тією самою частотою $\omega$, з~якою обертається диполь, і~амплітудою
	$I_0=2\pi p_m\omega a^2/(cR\ell^3)$. Так працюють генератори з~обертовим магнітом (розділ~\ref{ElMagInduction}) і~тахометри:
	частота струму в~нерухомій котушці дорівнює частоті обертання магніту.
\end{Example}



%% --------------------------------------------------------
\section{Енергія системи зв'язаних контурів}\hypertarget{coupled-circuits-energy}{}
%% --------------------------------------------------------

Формула~\eqref{eq:MagEnergyFormula} стосується \emph{одного} відокремленого контуру.
Виведемо тепер явний вираз для енергії системи довільної кількості $N$ магнітно зв'язаних
контурів зі струмами $I_1,\dots,I_N$, спираючись на щойно доведену теорему взаємності
$L_{ik}=L_{ki}$.

Як і~для одного контуру, приріст енергії дорівнює роботі джерел проти ЕРС індукції. У~$i$-му контурі за час $dt$ це робота $\dfrac{I_i}{c}\,d\Phi_i=\dfrac1{c^2}I_i\sum_kL_{ik}\,dI_k$ (контури нерухомі, $L_{ik}=\const$; ЕРС індукції дорівнює $-\frac1c\,d\Phi_i/dt$, потоки~--- за~\eqref{eq:FluxSumN}). Підсумувавши за~$i$, розіб'ємо приріст на~дві частини~--- від самоіндукції кожного контуру окремо та від взаємної індукції кожної пари контурів:
\begin{equation}\label{eq:EnergyIncrementDecomposition}
	dW = dW_\text{с}+dW_\text{в}, \qquad dW_\text{с}=\frac1{c^2}\sum_{i=1}^{N} L_iI_i\,dI_i, \qquad
	dW_\text{в}=\frac1{c^2}\sum_{\substack{i,k=1\\i\neq k}}^{N} L_{ik}I_i\,dI_k.
\end{equation}
Доданок $dW_\text{с}$ інтегрується так само, як і~для одного контуру
у~формулі~\eqref{eq:MagEnergyFormula}:
\begin{equation}\label{eq:SelfInductionEnergy}
	W_\text{с}=\frac1{2c^2}\sum_{i=1}^{N} L_iI_i^2.
\end{equation}
У~сумі $dW_\text{в}$ згрупуємо доданки попарно: для кожної пари індексів $i<k$ у~подвійній
сумі за $i\neq k$ трапляються рівно два доданки, $L_{ik}I_i\,dI_k$ і~$L_{ki}I_k\,dI_i$:
\begin{equation*}
	dW_\text{в} = \frac1{c^2}\sum_{i<k}\left(L_{ik}I_i\,dI_k+L_{ki}I_k\,dI_i\right).
\end{equation*}
Саме тут, і~лише тут, потрібна теорема взаємності: підставивши $L_{ki}=L_{ik}$, кожна
дужка перетворюється на повний диференціал добутку $I_iI_k$,
\begin{equation*}
	L_{ik}I_i\,dI_k+L_{ik}I_k\,dI_i = L_{ik}\,d(I_iI_k),
\end{equation*}
звідки
\begin{equation*}
	dW_\text{в} = \frac1{c^2}\sum_{i<k}L_{ik}\,d(I_iI_k) = \frac1{2c^2}\sum_{\substack{i,k=1\\i\neq k}}^{N}L_{ik}\,d(I_iI_k),
\end{equation*}
і, інтегруючи від нульових струмів,
\begin{equation}\label{eq:MutualEnergyFinal}
	W_\text{в}=\frac1{2c^2}\sum_{\substack{i,k=1\\i\neq k}}^{N}L_{ik}I_iI_k.
\end{equation}
Об'єднуючи~\eqref{eq:SelfInductionEnergy} і~\eqref{eq:MutualEnergyFinal}, отримуємо повну
енергію системи контурів:
\begin{equation}\label{eq:MultiLoopEnergy}
	\tcbhighmath{
		W = W_\text{с}+W_\text{в} = \frac1{2c^2}\sum_{i,k=1}^{N} L_{ik}I_iI_k = \frac1{2c}\sum_{i=1}^{N} I_i\Phi_i,
	}
\end{equation}
де в~останній рівності використано формулу \eqref{eq:FluxSumN}.

Для системи з~двох контурів формула~\eqref{eq:MultiLoopEnergy} набуває вигляду
\begin{equation}\label{eq:TwoLoopEnergy}
	W = \frac1{2c^2}\left(L_1I_1^2+L_2I_2^2+2L_{12}I_1I_2\right),
\end{equation}
де $L_1\equiv L_{11}$, $L_2\equiv L_{22}$~--- власні індуктивності контурів.



%% --------------------------------------------------------
\subsection*{Вираз енергії системи контурів через поля}
%% --------------------------------------------------------



Той самий результат~\eqref{eq:TwoLoopEnergy} можна отримати зовсім іншим шляхом~--- не
через роботу джерел, а~через \emph{локалізацію енергії в~полі}, узагальнюючи
формулу~\eqref{eq:MagEnergyDensityGeneral} на випадок поля, породженого одразу кількома
контурами (у~лінійному середовищі з~$\mu=\const$, як і~в~прикладі з~двома обмотками).
Нехай $\Bfield_1$~--- поле струму $I_1$, а~$\Bfield_2$~--- поле струму $I_2$; за принципом
суперпозиції $\Bfield=\Bfield_1+\Bfield_2$, і
\begin{equation*}
	W=\iiint\limits_V \frac{B^2}{8\pi\mu}\,dV = \iiint\limits_V \frac{B_1^2}{8\pi\mu}\,dV+\iiint\limits_V \frac{B_2^2}{8\pi\mu}\,dV+
	\iiint\limits_V \frac{\Bfield_1\cdot\Bfield_2}{4\pi\mu}\,dV.
\end{equation*}
Перші два інтеграли~--- це \enquote{власні} енергії кожного контуру окремо
($\tfrac12L_1I_1^2/c^2$ і~$\tfrac12L_2I_2^2/c^2$), а~третій, \enquote{перехресний},
доданок і~дає взаємну енергію $L_{12}I_1I_2/c^2$ з~формули~\eqref{eq:TwoLoopEnergy}; для
довільної кількості контурів міркування те саме, лише з~подвійною сумою перехресних
доданків замість одного.

Це трактування відразу дає кілька наслідків, не таких очевидних із формули роботи джерел:
\begin{enumerate}
	\item Енергія $W$~--- \emph{завжди додатна} за будь-яких струмів, оскільки під
	  інтегралом стоїть невід'ємна величина $B^2$.
	\item Енергія системи контурів \emph{не адитивна}: через перехресний доданок вона не
	  дорівнює сумі енергій окремих контурів.
	\item Рівність $L_{12}=L_{21}$ випливає із~симетрії перехресного доданка. Справді, $\Phi_{21}=L_{21}I_1/c$~--- потік поля $\Bfield_1$ крізь контур~2. Подавши $\Bfield_1=\rot\vect A_1$ і~проінтегрувавши частинами ($\rot\Hfield_2=\frac{4\pi}{c}\vect j_2$), маємо
	\begin{equation*}
		\iiint\limits_V \Bfield_1\cdot\Hfield_2\,dV=\iiint\limits_V \vect A_1\cdot\rot\Hfield_2\,dV
		=\frac{4\pi}{c}\,I_2\oint\limits_{C_2}\vect A_1\cdot d\vect l_2=\frac{4\pi}{c}\,I_2\Phi_{21},
	\end{equation*}
	тобто
	\begin{equation}\label{eq:MutualInductanceViaFieldIntegral}
		L_{21}=\frac{c^2}{4\pi I_1I_2}\iiint\limits_V \Bfield_1\cdot\Hfield_2\,dV.
	\end{equation}
	Але $\Bfield_1\cdot\Hfield_2=\mu\Hfield_1\cdot\Hfield_2=\Hfield_1\cdot\Bfield_2$, тож інтеграл не змінюється за перестановки індексів 1 і~2, і~$L_{12}=L_{21}$. Це друге доведення теореми взаємності, чинне й~для неоднорідного лінійного середовища $\mu=\mu(\vect r)$; \eqref{eq:MutualInductanceViaFieldIntegral}~--- друге означення взаємної індуктивності, через розподіл полів, а~не через потік крізь контур. В~однорідному середовищі підінтегральний вираз дорівнює $\Bfield_1\cdot\Bfield_2/\mu$.
\end{enumerate}
Отже, обидва підходи рівноправні й~дають той самий результат, що ще раз підтверджує:
магнітна енергія~--- це справді енергія \emph{поля}, розподілена в~просторі, а~не якась
умовна властивість самих контурів.

Та сама картина вже траплялась для електричного поля (розділ~\ref{ElEnergy}): розклад $E^2=E_1^2+E_2^2+2\Efield_1\cdot\Efield_2$ дає завжди додатну повну енергію~\eqref{eq:GeneralEnergyDensity} на відміну від знакозмінної взаємної енергії~\eqref{eq:SelfMutualEnergy}. Різниця лише в~тому, що власна енергія точкового заряду розбігається, а~власні енергії контурів $L_iI_i^2/2c^2$ скінченні й~додатні, тож знак може мати лише перехресний доданок $L_{ik}I_iI_k/c^2$.




%% --------------------------------------------------------
\subsection*{Обмеження на взаємну індуктивність}
%% --------------------------------------------------------



Вираз~\eqref{eq:TwoLoopEnergy} для енергії двох контурів~--- квадратична форма струмів
$I_1,I_2$. Енергія магнітного поля не може бути від'ємною за жодних значень струмів: $W=\iiint B^2/(8\pi\mu)\,dV\ge0$ (стан зі струмом не може мати меншу енергію, ніж стан без струму; див.~вище). Отже, ця квадратична форма повинна бути додатно визначеною. Квадратична форма $L_1I_1^2+2L_{12}I_1I_2+L_2I_2^2$ невід'ємна за всіх $I_1,I_2$ тоді й~лише тоді, коли її дискримінант недодатний, $L_{12}^2-L_1L_2\le0$, тобто
\begin{equation}\label{eq:InductanceInequality}
	\tcbhighmath{
		L_{12}^2 \le L_1L_2.
	}
\end{equation}
Відношення $k=L_{12}/\sqrt{L_1L_2}$ називають \emphx{коефіцієнтом зв'язку}; нерівність означає $|k|\le1$.

Фізично ця нерівність означає, що магнітний зв'язок між двома контурами не може бути
\enquote{сильнішим}, ніж дозволяють їхні власні індуктивності: частина енергії, яку
створює струм одного контуру, завжди залишається в~його \emph{власному} магнітному полі,
і~лише частина --- у~полі, зчепленому з~іншим контуром. Рівність $L_{12}^2=L_1L_2$ ($|k|=1$) означає
\emph{ідеальний} магнітний зв'язок (весь потік одного контуру повністю зчеплений з~іншим,
без розсіювання)~--- граничний, майже не досяжний на практиці випадок, до якого, однак,
впритул наближаються добре сконструйовані обмотки трансформатора на замкненому
феромагнітному осерді (\S~\ref{sec:IdealTransformer}).



%% --------------------------------------------------------
\section{Принцип роботи трансформаторів}\label{sec:IdealTransformer}\hypertarget{transformer-principle}{}
%% --------------------------------------------------------

Явище взаємної індукції лежить в~основі роботи \emphx{трансформаторів}~--- пристроїв для
підвищення чи зниження напруги змінного струму (рис.~\ref{pic:Transformer}). Первинна
й~вторинна обмотки з~$N_1$ і~$N_2$ витками укріплені на спільному замкненому
феромагнітному осерді; завдяки великому $\mu$ феромагнетика (розділ~\ref{MagFieldMedia})
магнітний потік майже повністю локалізований в~осерді й~пронизує обидві обмотки однаковим
потоком $\Phi$. Це~--- практична реалізація того самого \emph{ідеального} магнітного
зв'язку, граничного випадку $L_{12}^2=L_1L_2$ (формула~\eqref{eq:InductanceInequality}): розсіюванням
потоку поза осердям тут свідомо нехтують, а~отже, увесь потік, породжений первинною
обмоткою, зчеплений і з~вторинною. Змінна ЕРС $\emf_1$ у~первинній обмотці створює струм
$I_1$ і~змінний магнітний потік $\Phi$; цей потік пронизує й~витки вторинної обмотки,
наводячи в~ній ЕРС взаємної індукції, а~в~первинній~--- ЕРС самоіндукції.

\begin{figure}[h!]\centering
	\localinput{Transformer.tikz}
	\caption{Принципова схема трансформатора\label{pic:Transformer}}
\end{figure}

Закон Ома для первинного кола: $\emf_1 - \frac1c\dfrac{d}{dt}(N_1\Phi) = I_1R_1$. У~режимі
швидкозмінних полів (типовому для трансформаторів) $I_1R_1 \ll \emf_1,\
\tfrac{d}{dt}(N_1\Phi)$, тому у~цьому наближенні
\begin{equation}\label{eq:TransformerEMF1}
	\emf_1 \approx \frac{N_1}c\frac{d\Phi}{dt},
\end{equation}
а~ЕРС у~вторинній обмотці, оскільки потік $\Phi$ в~осерді той самий,
\begin{equation}\label{eq:TransformerEMF2}
	\emf_2 = -\frac{N_2}c\frac{d\Phi}{dt} = -\frac{N_2}{N_1}\emf_1
\end{equation}
(знак \enquote{мінус} означає лише, що за однакового напрямку обходу обмоток ЕРС вторинної обмотки, як і~ЕРС самоіндукції первинної, протилежна за знаком ЕРС джерела; знак напруги на виході залежить від напрямку намотування). Відношення
\begin{equation}\label{eq:TransformationRatio}
	\tcbhighmath{
		\eta = \frac{N_2}{N_1}
	}
\end{equation}
називають \emphx{коефіцієнтом трансформації}: якщо $\eta>1$~--- трансформатор
\emphx{підвищувальний} (збільшує ЕРС), якщо $\eta<1$ --- \emphx{знижувальний} (зменшує
ЕРС).

А~як перетворюється струм? Тут працює закон збереження енергії. У~наближенні ідеального трансформатора~--- без омічних втрат ($R_1,R_2\to0$) і~без розсіювання потоку~--- в~обмотках не накопичується енергія понад ту, що в~кожен момент проходить крізь них: уся потужність $I_1\emf_1$, яку джерело вкладає в~первинне коло, через спільне поле передається навантаженню вторинного кола, де її віддає ЕРС індукції $\emf_2$,
\begin{equation*}
	I_1\emf_1 = I_2\emf_2
\end{equation*}
(струми й~ЕРС відліковано в~одному й~тому самому напрямку обходу кожної обмотки). Підставляючи~\eqref{eq:TransformerEMF2}, знаходимо коефіцієнт трансформації \emph{струмів}:
\begin{equation}\label{eq:TransformerCurrentRatio}
	\tcbhighmath{
		\frac{I_2}{I_1} = \frac{\emf_1}{\emf_2} = -\frac1\eta.
	}
\end{equation}
Знак \enquote{мінус} тут невипадковий: за однакового обходу струми в~обмотках течуть протилежно, ампер-витки компенсуються, $N_1I_1+N_2I_2=0$. Тому струм вторинної обмотки не підсилює, а~гасить потік первинного; потік $\Phi$ в~осерді задає лише ЕРС~$\emf_1$, див.~\eqref{eq:TransformerEMF1}.

Той самий результат виходить і~без енергії. За теоремою про циркуляцію $\oint\Hfield\cdot d\vect l=\frac{4\pi}{c}(N_1I_1+N_2I_2)$, а~в~осерді $H=B/\mu\to0$ при $\mu\to\infty$, тож $N_1I_1+N_2I_2\to0$. Звідси й~відповідь на питання, чому за розімкненої вторинної обмотки ($I_2=0$) трансформатор не споживає струму: тече лише намагнічувальний струм, що створює потік $\Phi$ в~осерді, і~він зменшується як $1/\mu$. Нарешті, за ідеального зв'язку $L_{12}=\sqrt{L_1L_2}$, а~$L_1\propto N_1^2$, $L_2\propto N_2^2$, тож $\eta=N_2/N_1=\sqrt{L_2/L_1}$.

Отже, за модулем у~скільки разів трансформатор підвищує напругу, у~стільки ж разів він знижує
струм,~--- і~навпаки: закон збереження енергії не дозволяє одночасно виграти і~в~напрузі,
і~в~струмі. Саме тому передавати електроенергію на великі відстані вигідно при високій
напрузі й, відповідно, малому струмі (а~отже, і~з~малими тепловими втратами $I^2R$
у~дротах лінії передачі), а~споживачам подавати вже при низькій, безпечній напрузі~---
і~саме взаємна індукція, а~не який-небудь прямий електричний контакт, дозволяє так легко
узгоджувати ці дві вимоги.

Тут ми розглянули лише ідеалізований трансформатор~--- без омічних втрат в~обмотках і~без
розсіювання потоку,~--- щоб якнайпростіше дістати сам факт існування коефіцієнта
трансформації з~енергетичних міркувань. Реальний трансформатор як коло змінного струму, з
вхідним імпедансом, внесеним опором, режимами холостого ходу й~короткого замикання,
втратами в~міді та сталі, ККД, а~також практичними застосуваннями (узгодження
навантаження, автотрансформатори, трифазні трансформатори, передача електроенергії на
великі відстані) детально розглянуто в~\S~\ref{sec:transformer}, присвяченому змінному
струму.



%% --------------------------------------------------------
\section{Магнітна енергія струмів як функція стану}\label{sec:CurrentPotentialEnergy}\hypertarget{current-potential-energy}{}
%% --------------------------------------------------------

Вираз~\eqref{eq:MultiLoopEnergy} для енергії системи контурів за традицією називають \emphx{потенціальною енергією струмів}. Назва потребує застереження. Магнітне поле \emph{не є потенціальним}: там, де течуть струми, $\rot\Bfield\neq0$, тож $\Bfield\neq-\grad\phi$ в~усьому просторі, і~$W$ не є потенціалом цього поля. Це \emph{функція стану} системи контурів, $W(I_1,\dots,I_N)$ за~заданих $L_{ik}$: вона не залежить від того, в~якому порядку й~як швидко встановлювалися струми, а~механічні сили виходять із~неї диференціюванням (\S~\ref{sec:MagPonderomotor}).

Строго кажучи, в~термінах аналітичної механіки $W$~--- це \enquote{магнітна кінетична енергія} системи: квадратична форма струмів (аналогів швидкостей) з~коефіцієнтами $L_{ik}$ у~ролі узагальнених мас. Тому знак сили інший, ніж для механічної потенціальної енергії. За сталих струмів $F_x=+(\partial W/\partial x)_I$, тобто роль потенціальної енергії відіграє $-W$ (рівновага стійка в~максимумі $W$). За сталих потоків $F_x=-(\partial W/\partial x)_\Phi$, і~тоді потенціальною справді є $W=\Phi^2/2L$ (рівновага в~мінімумі).



%% --------------------------------------------------------
\subsection*{Вираз через векторний потенціал}
%% --------------------------------------------------------



Тепер покажемо, як ту саму величину~\eqref{eq:MultiLoopEnergy} виразити безпосередньо
через векторний потенціал поля, не вдаючись до обчислення геометричних індуктивностей
$L_{ik}$. Для лінійного середовища ($\Bfield=\mu\Hfield$, $\mu=\const$) скористаємось
означенням $\Bfield=\rot\vect A$ і~тим, що для системи лінійних струмів векторний
потенціал дорівнює сумі внесків окремих контурів, $\vect A=\sum_i\vect A_i$, де
\begin{equation}\label{eq:VectorPotentialLineCurrentSum}
	\vect A_i(\vect r) = \frac{\mu I_i}{c}\oint\limits_{C_i} \frac{d\vect l_i}{|\vect r-\vect r_i|}.
\end{equation}
Підставляючи $\Bfield=\rot\vect A$ у~вираз для повної енергії поля і~інтегруючи частинами
(з~урахуванням того, що на нескінченності поля спадають достатньо швидко), отримуємо:
\begin{equation}\label{eq:MagEnergyViaVectorPotential}
	W = \frac1{8\pi}\iiint\limits_V \Bfield\cdot\Hfield\,dV
	= \frac1{8\pi}\iiint\limits_V \rot\vect A\cdot\Hfield\,dV
	= \frac1{8\pi}\iiint\limits_V \vect A\cdot\rot\Hfield\,dV
	= \frac1{2c}\iiint\limits_V \vect A\cdot\vect j\,dV,
\end{equation}
де на останньому кроці використано рівняння $\rot\Hfield=\frac{4\pi}{c}\vect j$ ($\vect j$~--- густина струму вільних носіїв). Для
системи тонких контурів $\vect j\,dV\to I_i\,d\vect l_i$,
і~формула~\eqref{eq:MagEnergyViaVectorPotential} переходить у
\begin{equation}\label{eq:MagEnergyLineCurrents}
	W = \frac1{2c}\sum_{i=1}^{N} I_i\oint\limits_{C_i}\vect A\cdot d\vect l_i
	= \frac1{2c}\sum_{i=1}^{N} I_i\Phi_i,
\end{equation}
що точно збігається з~отриманим раніше виразом~\eqref{eq:MultiLoopEnergy}. Отже, обидва
означення~--- через контурні індуктивності $L_{ik}$ і~через суцільний розподіл векторного
потенціалу $\vect A$~--- описують ту саму фізичну величину. Для нескінченно тонкої нитки $\vect A$ на самій нитці (а~отже й~власна енергія) розбігається логарифмічно, тож дріт вважаємо скінченної товщини.


Зіставимо магнітний вираз~\eqref{eq:MagEnergyViaVectorPotential} із його електростатичним
аналогом, де енергія системи зарядів записується через скалярний потенціал $\phi$:
\begin{equation*}
	W_e = \frac12\iiint\limits_V \rho\,\phi\,dV
	= \frac1{8\pi}\iiint\limits_V \Efield\cdot\Dfield\,dV.
\end{equation*}
Ми бачимо повну структурну відповідність: енергія системи джерел~--- це \emph{інтеграл від добутку джерела на його потенціал} ($\rho\phi$ в~електриці, $\vect A\cdot\vect j$ у~магнетизмі). Магнітна формула зручна для розподілених струмів (об'ємних, поверхневих), де поняття індуктивності контуру втрачає сенс. Потенціали визначені з~точністю до калібрування, і~магнітний інтеграл від неї не залежить: заміна $\vect A\to\vect A+\grad\chi$ додає $\iiint\vect j\cdot\grad\chi\,dV=-\iiint\chi\,\divg\vect j\,dV=0$ (струми сталі, $\divg\vect j=0$, і~займають скінченну область). В~електростатиці ж заміна $\phi\to\phi+\const$ змінює $\iiint\rho\phi\,dV$ на~$\const\cdot Q$, тому там фіксують $\phi(\infty)=0$.


\begin{Example}{Енергія соленоїда через векторний потенціал}

	Застосуємо нову формулу~\eqref{eq:MagEnergyViaVectorPotential} до довгого соленоїда
	радіуса $R$, довжини $\ell$ з~густиною намотування $n$ витків на одиницю довжини, яким
	тече струм $I$.

	Поле всередині соленоїда однорідне, $B=4\pi\mu nI/c$. Векторний потенціал
	у~циліндричних координатах, що задовольняє умову $\Bfield=\rot\vect A$, дорівнює:
	\begin{equation}\label{eq:SolenoidVectorPotential}
		A_\phi(r) =
		\begin{cases}
			\dfrac{B r}{2},    & r < R, \\[10pt]
			\dfrac{B R^2}{2r}, & r > R.
		\end{cases}
	\end{equation}
	Струм намотки зосереджений на поверхні $r=R$, де $\vect j\,dV = I\,n\,\ell\,d\vect l$.
	Підставляючи $A_\phi(R) = BR/2$ у~формулу~\eqref{eq:MagEnergyViaVectorPotential}:
	\begin{equation*}
		W = \frac1{2c}\iiint\limits_V \vect A\cdot\vect j\,dV
		= \frac1{2c}\cdot I\cdot N\cdot A_\phi(R)\cdot 2\pi R
		= \frac1{2c}\cdot I\cdot N\cdot \frac{BR}{2}\cdot 2\pi R
		= \frac{\pi N I B R^2}{2c}.
	\end{equation*}
	Підставляючи $B=4\pi\mu nI/c$ та $N=n\ell$, маємо:
	\begin{equation*}
		W = \frac{4\pi\mu N^2 S}{\ell}\cdot\frac{I^2}{2c^2} = \frac{LI^2}{2c^2},
	\end{equation*}
	де $L=4\pi\mu N^2S/\ell$~--- індуктивність соленоїда~\eqref{eq:SolenoidInductance}.
	Результат точно збігається з~раніше отриманою формулою, але розрахований безпосередньо
	з~характеристик $\vect A$ та $\vect j$.
\end{Example}



%% --------------------------------------------------------
\section{Пондеромоторні сили в~магнітному полі}\label{sec:MagPonderomotor}\hypertarget{magnetic-ponderomotive-forces}{}
%% --------------------------------------------------------



%% --------------------------------------------------------
\subsection*{Метод віртуальних переміщень для магнітних сил}
%% --------------------------------------------------------



Повернімося\index{Віртуальних переміщень метод!для магнітних сил} тепер до величини $W$, яку ми назвали потенціальною енергією струмів. Так
само, як і~в~електростатиці (розділ~\ref{ElEnergy}), силу, що діє на провідник чи магнетик
з~боку магнітного поля, зручно шукати методом віртуальних переміщень, спираючись на закон
збереження енергії. Робота джерел струму $\delta A_\text{дж}$ витрачається на теплоту
$\delta Q$, на приріст магнітної енергії $dW$ і~на механічну роботу $\delta A_\text{мех}$:
\begin{equation}\label{eq:PonderomotiveWork}
	\delta A_\text{дж} = \delta Q + dW + \delta A_\text{мех}.
\end{equation}
З~іншого боку, за законом Ома, робота джерела дорівнює $\delta A_\text{дж} = \delta Q -
\emf_\text{сі}I\, dt$ (ЕРС самоіндукції~--- частина роботи джерела, витрачена на подолання
цієї ЕРС). Порівнюючи два вирази для $\delta A_\text{дж}$:
\begin{equation*}
	-\emf_\text{сі}I\, dt = dW + \delta A_\text{мех}.
\end{equation*}

\emph{Випадок $I=\const$} (струм у~колі підтримується сталим джерелом). Тоді, за
формулою~\eqref{eq:SelfInductionEMF}, $-\emf_\text{сі}I\,dt = \dfrac{I^2}{c^2}dL$, тоді як
приріст магнітної енергії за формулою~\eqref{eq:MagEnergyFormula} становить лише половину
цього, $dW = \dfrac{I^2}{2c^2}dL$ (оскільки $I=\const$, а~не $\Phi$). Отже, $\delta
A_\text{мех} = +dW$, і~сила
\begin{equation}\label{eq:ForceConstI}
	\tcbhighmath{
		F_x = +\left(\frac{dW}{dx}\right)_I.
	}
\end{equation}

\emph{Випадок $\Phi=\const$} (наприклад, контур з~надпровідника). Тоді $\emf_\text{сі} =
0$ (адже $d\Phi/dt=0$), тому $\delta A_\text{мех} = -dW$, і
\begin{equation}\label{eq:ForceConstPhiM}
	\tcbhighmath{
		F_x = -\left(\frac{dW}{dx}\right)_\Phi.
	}
\end{equation}

Знаки формул~\eqref{eq:ForceConstI} і~\eqref{eq:ForceConstPhiM} такі самі, як і~у~відповідних формулах
електростатики~\eqref{eq:ForceConstQ}, \eqref{eq:ForceConstPhiE} (розділ~\ref{ElEnergy}), якщо правильно зіставити режими.
Сталий струм відповідає сталій різниці потенціалів: в~обох випадках джерело підтримує режим і~виконує роботу $2\,dW$, половина
якої йде на приріст енергії поля, а~половина~--- на механічну роботу, тому $F_x=+\partial W/\partial x$. Сталий потік відповідає
сталому заряду: система ізольована від джерел, і~механічна робота виконується за рахунок енергії поля,
$F_x=-\partial W/\partial x$. Плутанина виникає, якщо формально зіставляти струм із зарядом (енергія $LI^2/2c^2$ зовні схожа на
$q^2/2C$),~--- тому варто щоразу з'ясовувати, хто виконує роботу.

Перевіримо знак на відомому факті. Два співвісні кільця з~однаково напрямленими струмами: потоки додаються, $L_{12}>0$, і~$L_{12}$ зростає при зближенні. Перехресна частина енергії дорівнює $L_{12}I_1I_2/c^2$, тож за $I=\const$ сила $F_x=+\dfrac{I_1I_2}{c^2}\dfrac{\partial L_{12}}{\partial x}<0$, де $x$~--- відстань між кільцями: вона напрямлена до зближення. Це відоме притягання співнапрямлених струмів (розділ~\ref{MagFieldVacuum}); знак \enquote{плюс} перед похідною дає саме правильний результат.

\begin{Attention}
	Строго кажучи, роботу сил поля за сталої температури визначає не сама енергія $W$,
	а~\emphx{вільна енергія} $\mathcal F$ (термодинамічний потенціал, що визначає роботу за $T=\const$). Різниця між ними виникає лише тоді, коли властивості середовища
	(наприклад, магнітна проникність $\mu$) залежать від температури й~віртуальне
	переміщення супроводжується тепловими ефектами. У~прикладах цього розділу
	$\mu=\const$, тому $\mathcal F$ і~$W$~--- буквально та сама величина,
	і~формули~\eqref{eq:ForceConstI}, \eqref{eq:ForceConstPhiM} застосовуються без змін.
\end{Attention}

\begin{Example}{Втягування стержня з~магнетика в~соленоїд}

	Знайдемо, користуючись загальним методом, силу, з~якою втягується у~вакуумний соленоїд
	довгий циліндричний стержень (рис.~\ref{pic:RodInSolenoid}) з~магнітною проникністю $\mu$ і~площею поперечного
	перерізу $S$. По соленоїду, що має густину намотування $n=N/\ell$, циркулює струм $I$.

	Вважатимемо, що струм у~соленоїді підтримується сталим, тому користуємось формулою при
	сталому струмі,
	\begin{equation}\label{eq:WorkConstCurrentRod}
		\delta A_\text{поле} = f\,\delta x = \left(dW\right)_{I} = \frac{I}{2c}\,d\Phi
	\end{equation}
	(тут $\Phi$~--- повне потокозчеплення соленоїда, тобто сума потоків крізь усі його витки, і~$W=I\Phi/2c$, як
	у~формулі~\eqref{eq:MagEnergyFormula}). Далі, оскільки соленоїд і~стержень довгі, а~розглядаємо ділянку далеко від їхніх кінців, за теоремою про циркуляцію $H=4\pi n I/c$ і~всередині стержня, і~поза ним (дотична складова $\Hfield$ неперервна на бічній поверхні стержня, а~розмагнічувальними ефектами на кінцях нехтуємо; розділ~\ref{MagFieldMedia}), звідки $I=cH/4\pi n$, а~отже
	\begin{equation}\label{eq:FreeEnergyDiffViaH}
		\left(dW\right)_{I} = \frac1{8\pi n}H\,d\Phi.
	\end{equation}

	Знайдемо зміну потоку $d\Phi$ при зміщенні стержня на $\delta x$. Потокозчеплення $\Phi$
	змінюється за рахунок того, що частина довжини соленоїда, яка досі
	містила вакуум ($\mu=1$), заповнюється речовиною з~проникністю $\mu$: при зміщенні
	стержня на $\delta x$ число витків, що охоплюють вакуумну частину, зменшується на
	$n\,\delta x$, а~число витків, що охоплюють частину з~магнетиком, на стільки ж
	збільшується. Тому, нехтуючи неоднорідністю поля поблизу межі розділу й~кінців
	соленоїда (об'єм цих ділянок при малому $\delta x$ практично не змінюється)
	і~використовуючи $B_\text{стержень}=\mu H$, $B_\text{вакуум}=H$
	(розділ~\ref{MagFieldMedia}):
	\begin{equation}\label{eq:FluxChangeRod}
		d\Phi = n\,\delta x\left(B_\text{стержень}-B_\text{вакуум}\right)S = n\,\delta x\,(\mu-1)HS.
	\end{equation}
	Підставляючи~\eqref{eq:FluxChangeRod} у~\eqref{eq:FreeEnergyDiffViaH}, знаходимо
	\begin{equation*}
		\left(dW\right)_{I} = \frac{\mu-1}{8\pi}H^2S\,\delta x,
	\end{equation*}
	звідки, за формулою~\eqref{eq:WorkConstCurrentRod},
	\begin{equation}\label{eq:ForceOnRod}
		f = \frac{\left(dW\right)_{I}}{\delta x} = \frac{\mu-1}{8\pi}H^2S.
	\end{equation}
	При $\mu>1$ (парамагнетик) сила додатна~--- стержень втягується в~соленоїд; при
	$\mu<1$ (діамагнетик) сила від'ємна~--- стержень виштовхується з~нього.
\end{Example}

\begin{figure}[h!]\centering
	\includegraphics[width=0.55\linewidth]{pull_in_solenoid.jpg}
	\caption{Осердя втягується в~котушку зі~струмом~--- приклад пондеромоторної сили в~магнітному полі\label{pic:RodInSolenoid}}
\end{figure}

Формула~\eqref{eq:ForceOnRod}~--- ще одне підтвердження загального правила: парамагнетики
($\mu>1$) втягуються в~область сильнішого поля, а~діамагнетики ($\mu<1$)~--- виштовхуються
з~неї. Сила пропорційна $H^2$ (тобто не залежить від напрямку поля) і~площі перерізу
стержня $S$~--- чим \enquote{товщий} стержень, тим більша втягувальна сила. Порівняйте з~формулою~\eqref{eq:ElectromagnetForce} для підйомної сили електромагніта: і~там, і~тут сила дорівнює добутку площі на~величину порядку густини енергії поля. Енергетичний метод дає змогу знаходити силу в~найрізноманітніших конфігураціях, не вдаючись щоразу до прямого обчислення сил Ампера.

\begin{Example}{Підйомна сила електромагніта}
	Розглянемо електромагніт~--- П-по\-діб\-не осердя з~феромагнетика проникності $\mu$,
	обмотане $N$ витками зі струмом $I$, до полюсів якого прикладено якір (пластину),
	відокремлений повітряним зазором завтовшки $x$ з~кожного боку. Знайдемо силу
	притягання якоря.

	За теоремою про циркуляцію для $\Hfield$, застосованою до контуру, що проходить через
	осердя (довжиною $\ell$) і~двічі перетинає зазор (загальною товщиною $2x$):
	\begin{equation*}
		\frac{B}{\mu}\ell + B\cdot 2x = \frac{4\pi}{c}IN \quad\Rightarrow\quad B = \frac{4\pi\mu}{c}\,\frac{NI}{\ell+2\mu x}
	\end{equation*}
	(у~повітряному зазорі $\mu_\text{пов}\approx1$, тому там $\Hfield=\Bfield$; поле
	в~зазорі й~осерді те саме завдяки неперервності $B_n$ на межі,
	розділ~\ref{MagFieldMedia}); розсіюванням потоку й~крайовими ефектами нехтуємо. Магнітний потік крізь переріз $S$ магніту: $\Phi = NBS =
	\dfrac{4\pi\mu}{c}\dfrac{N^2I}{\ell+2\mu x}S$.

	За умови сталого струму ($I=\const$) зміна енергії поля при віртуальній зміні зазору
	$dW = \dfrac{I\, d\Phi}{2c}$ (з~формули~\eqref{eq:MagEnergyFormula}, $W=I\Phi/(2c)$
	при $I=\const$), і, за формулою~\eqref{eq:ForceConstI},
	\begin{equation*}
		F_x = \left(\frac{dW}{dx}\right)_I = \frac{I}{2c}\left(\frac{d\Phi}{dx}\right)_I = -\frac{4\pi\mu^2}{c^2}\,\frac{N^2I^2}{(\ell+2\mu x)^2}S =
		-\frac{B^2}{8\pi}\cdot 2S.
	\end{equation*}
	Отже,
	\begin{equation}\label{eq:ElectromagnetForce}
		F_x = -\frac{B^2}{4\pi}S.
	\end{equation}
	Від'ємний знак означає, що сила~--- сила \emph{притягання}: якір втягується в~зазор,
	зменшуючи $x$,~--- цілком у~дусі загального висновку про те, що пара- й~феромагнетики (тут~---
	феромагнітний якір) завжди втягуються в~область сильнішого поля
	(формула~\eqref{eq:ForceOnRod}).
\end{Example}

\begin{figure}[h!]\centering
	\localinput{electromagnet.tikz}
	\caption{Електромагніт з~якорем, відокремленим повітряним зазором $x$\label{pic:ElectromagnetForce}}
\end{figure}



%% --------------------------------------------------------
\section*{Підсумки}
\phantomsection \addcontentsline{toc}{section}{Підсумки}
%% --------------------------------------------------------

Шлях цього розділу був обернений до того, яким вводилась енергія електричного поля: там
вона народжувалась \emph{напряму} з~роботи проти сили Кулона. Тут так не вийшло: магнітна
сила, за Лоренцем, роботи над зарядами не виконує ніколи~--- ні для сталого поля, ні для
змінного. Але, змінюючись, магнітне поле породжує інше поле~--- \emphx{вихрове
електричне}~--- і~вже воно виконує роботу над зарядами. Саме цю роботу, на подолання ЕРС
\emphx{самоіндукції}, і~виконує джерело струму, а~не роботу проти сили Ампера.

З~цієї роботи ми й~отримали густину енергії $w=BH/8\pi$~--- за формою тотожну
електростатичній $w=ED/8\pi$ (розділ~\ref{ElEnergy}). Та сама логіка лежить і~в~основі
пондеромоторних сил (втягування магнетика в~соленоїд, притягання якоря електромагніта). На
макрорівні здається, що механічну роботу виконує саме магнітна сила. Проте для заряду, що
рухається \emph{разом} із провідником, ця ж сила одночасно гальмує дрейфовий струм.
Робота магнітної сили над макроскопічним рухом тіла в~точності компенсується її ж
від'ємною роботою над струмом~--- так що сумарно, як і~належить, сила Лоренца роботи не
виконує. У~мові ж електродинамічних контурів рух тіла змінює індуктивність $L(x)$,
викликаючи додаткову ЕРС індукції. Саме джерело струму, долаючи цю ЕРС, компенсує
гальмування носіїв і~здійснює підсумкову роботу, яка за сталого струму ділиться порівну: одна її частина йде
на поповнення енергії поля з~густиною $w = BH/8\pi$, а~інша~--- переходить у~механічний
розгін тіла.

Решта розділу~--- наслідки з~цієї енергетики. Взаємна індукція описується коефіцієнтами $L_{ik}$, які за теоремою взаємності симетричні, $L_{12}=L_{21}$. Енергія системи контурів $W=\frac1{2c^2}\sum L_{ik}I_iI_k$ додатна, звідки $L_{12}^2\le L_1L_2$; граничний випадок ідеального зв'язку ($|k|=1$) реалізується в~трансформаторі, де $\emf_2/\emf_1=-N_2/N_1$, а~струми змінюються в~протилежний бік. Сили знаходимо з~$W$ методом віртуальних переміщень: $F_x=+(\partial W/\partial x)_I$ за сталих струмів і~$F_x=-(\partial W/\partial x)_\Phi$ за сталих потоків.

Звідси головний висновок розділу: магнітне поле~--- не \emph{виконавець},
а~\emph{посередник}. Енергію воно \emph{зберігає}, але обмінюється нею із зарядами лише через електричне поле: вихрове, яке воно породжує, змінюючись у~часі, а~для провідника, що рухається,~--- поле $\Efield'=\vect v\times\Bfield/c$ в~його власній системі відліку.



%% --------------------------------------------------------
\section*{Задачі}



\phantomsection \addcontentsline{toc}{section}{Задачі}
%% --------------------------------------------------------

% задачі цього розділу короткі: не розриваємо їх між сторінками

\begin{problem}
	Замкнений плоский контур із постійним струмом \(I\) поміщено в~зовнішнє однорідне
	магнітне поле \(\Bfield\). Поле повертає контур, орієнтуючи його магнітний момент
	уздовж \(\Bfield\). Відомо, що сила Лоренца не виконує роботи над носіями заряду.
	Поясніть, за рахунок якого джерела і~через який мікроскопічний механізм здійснюється
	механічний розгін контуру. Запишіть енергетичний баланс. Чому дорівнює робота джерела, що підтримує струм?
	\begin{solution}
		Магнітний момент $p_m=IS/c$, механічна робота повороту $A=-\Delta U=\Delta(\vect p_m\cdot\Bfield)=I\,\Delta\Phi/c$, де $\Phi=BS\cos\theta$~--- потік зовнішнього поля крізь контур. Під час повороту ЕРС індукції $-\frac1c\,d\Phi/dt$ гальмує струм, і~джерело, що підтримує $I=\const$, виконує роботу $I\,d\Phi/c$, яка дорівнює механічній (зміну енергії зовнішнього поля оплачують його власні джерела; нагрівання не враховуємо). Сила Лоренца роботи не виконує: її складові, пов'язані з~рухом контуру й~з~дрейфом носіїв, компенсують одна одну.
	\end{solution}
\end{problem}

\begin{problem}
	У~те саме однорідне магнітне поле \(\Bfield\) поміщено намагнічену феромагнітну
	стрілку з~таким самим за величиною магнітним моментом \(\vect{p}_m\). Вона
	повертається вздовж поля, не маючи макроскопічного струму чи джерела живлення.
	Спираючись на формулу потенціальної енергії магнітного диполя у~зовнішньому
	полі~\eqref{eq:PotentialEnergyDipole}, поясніть, як механічна робота повороту
	пов'язана зі зміною цієї енергії, і~які саме внутрішні сили безпосередньо повертають
	стрілку. Запишіть енергетичний баланс та порівняйте його із балансом для контуру зі
	струмом.
\end{problem}

\begin{problem}\label{prb:MRIEnergy}
	\emph{Енергія магніту томографа.} Надпровідний магніт томографа створює поле $B=3$~Тл ($3\cdot10^4$~Гс) в~об'ємі близько
	$1$~м$^3$. Оцініть енергію, запасену в~полі, і~порівняйте з~енергією вибуху $1$~кг тротилу ($4{,}2$~МДж). Чому аварійна втрата
	надпровідності (\enquote{квенч}) такого магніту небезпечна?
	\begin{solution}
		$w=\dfrac{B^2}{8\pi}\approx3{,}6\cdot10^{7}$~ерг/см$^3$, $W=wV\approx3{,}6\cdot10^{13}$~ерг $\approx3{,}6$~МДж~--- майже як $1$~кг тротилу. При квенчі ця енергія за секунди виділяється як тепло в~обмотці й~випаровує рідкий гелій.
	\end{solution}
\end{problem}

\begin{problem}\label{prb:CoaxInductance}
	\emph{Індуктивність коаксіального кабелю.} По внутрішньому тонкостінному циліндру радіуса $a$ коаксіального кабелю тече
	струм $I$, по зовнішньому радіуса $b$~--- зворотний струм. Знайдіть енергію магнітного поля на одиницю довжини кабелю й~звідси
	індуктивність одиниці довжини. Обчисліть її для $a=0{,}5$~мм, $b=3{,}5$~мм ($1$~Гн $=10^9$~см).
	\begin{solution}
		поле є лише між циліндрами, $H=\dfrac{2I}{cr}$; $\dfrac{W}{\ell}=\displaystyle\int_a^b\frac{H^2}{8\pi}2\pi r\,dr=\frac{I^2}{c^2}\ln\frac ba$, тому $\dfrac{L}{\ell}=2\ln\dfrac ba\approx3{,}9$ (см на см) $\approx0{,}39$~мкГн/м.
	\end{solution}
\end{problem}

\begin{problem}\label{prb:CoaxialRingsNumbers}
	\emph{Взаємна індуктивність кілець.} Для двох концентричних кілець в~одній площині з~радіусами $R_1=10$~см і~$R_2=1$~см
	(рис.~\ref{pic:CoaxialLoops}) знайдіть взаємну індуктивність у~сантиметрах і~в~генрі. Яка ЕРС наводиться в~малому кільці,
	якщо струм у~великому змінюється зі швидкістю $10^3$~А/с?
	\begin{solution}
		$L_{12}=\dfrac{c\Phi_{12}}{I}=\dfrac{2\pi^2R_2^2}{R_1}\approx2$~см $=2$~нГн; $|\emf|=L_{12}\left|\dfrac{dI}{dt}\right|$ (у~СІ) $\approx2$~мкВ.
	\end{solution}
\end{problem}

\begin{problem}\label{prb:ElectromagnetNumbers}
	\emph{Підйомна сила електромагніта.} Індукція в~зазорі електромагніта (рис.~\ref{pic:ElectromagnetForce}) $B=10$~кГс, площа
	кожного з~двох полюсів $S=10$~см$^2$. Знайдіть силу притягання якоря.
	\begin{solution}
		за~\eqref{eq:ElectromagnetForce} $F=\dfrac{B^2S}{4\pi}\approx8\cdot10^{7}$~дин $=800$~Н (близько $80$~кгс).
	\end{solution}
\end{problem}

\begin{problem}\label{prb:MagneticPressure}
	\emph{Магнітний тиск.} Методом віртуальних переміщень за сталого потоку покажіть, що на обмотку довгого соленоїда з~полем $B$
	діє тиск $p=B^2/8\pi$, напрямлений назовні. Оцініть його для $B=10$~Тл.
	\begin{solution}
		$W=\dfrac{B^2}{8\pi}S\ell=\dfrac{\Phi^2\ell}{8\pi S}$ при $\Phi=BS=\const$; при збільшенні площі перерізу $dW=-\dfrac{B^2}{8\pi}\ell\,dS=-p\,dV$, тож $p=\dfrac{B^2}{8\pi}$. Для $B=10^5$~Гс $p\approx4\cdot10^{8}$~дин/см$^2\approx4\cdot10^{7}$~Па, тобто близько $400$~атм.
	\end{solution}
\end{problem}

\begin{problem}\label{prb:SeriesCoils}
	\emph{Послідовне з'єднання котушок.} Дві котушки з~індуктивностями $L_1$, $L_2$ і~взаємною індуктивністю $L_{12}$ з'єднано послідовно. Знайдіть повну індуктивність за узгодженого (потоки котушок додаються) і~зустрічного (потоки віднімаються) вмикання. Чому нерівність $L_{12}^2\le L_1L_2$ гарантує $L\ge0$? Обчисліть для $L_1=L_2=10$~мГн і~коефіцієнта зв'язку $k=0{,}5$.
	\begin{solution}
		той самий струм $I$ в~обох: $W=\dfrac{1}{2c^2}\left(L_1+L_2\pm2L_{12}\right)I^2$, тому $L=L_1+L_2\pm2L_{12}$. Найменше значення $L_1+L_2-2\sqrt{L_1L_2}=(\sqrt{L_1}-\sqrt{L_2})^2\ge0$. Для $L_{12}=kL=5$~мГн: $30$~мГн (узгоджено) і~$10$~мГн (зустрічно).
	\end{solution}
\end{problem}

\begin{problem}\label{prb:TransformerNumbers}
	\emph{Знижувальний трансформатор.} Первинну обмотку з~$N_1=2200$ витків ввімкнено в~мережу $220$~В, до вторинної з~$N_2=110$ витків підключено навантаження $11$~Ом. Вважаючи трансформатор ідеальним, знайдіть напругу й~струм у~вторинному колі та струм у~первинному. Електроенергію потужністю $1$~МВт передають лінією опором $10$~Ом; порівняйте втрати $I^2r$ за напруги $10$~кВ і~$100$~кВ.
	\begin{solution}
		$\eta=N_2/N_1=0{,}05$; $U_2=\eta U_1=11$~В, $I_2=U_2/R=1$~А, $I_1=\eta I_2=0{,}05$~А; $P=U_2I_2=U_1I_1=11$~Вт. У~лінії $I=P/U$ дорівнює $100$~А і~$10$~А, втрати $I^2r=100$~кВт ($10\%$) і~$1$~кВт ($0{,}1\%$): у~$100$ разів менше.
	\end{solution}
\end{problem}

\begin{Questions}
	\begin{quizlist}
		\item Виведіть енергію магнітного поля контуру зі струмом $W = LI^2/2c^2$ з~роботи
		  джерела ЕРС. На що йде решта роботи джерела?
		\item Як звідси отримати густину енергії магнітного поля $w = BH/8\pi$ для
		  соленоїда? Як це виведення узагальнюється на довільне поле і~яку роль у~ньому
		  відіграє вихрове електричне поле?
		\item Дайте означення коефіцієнтів взаємної індуктивності $L_{12}$, $L_{21}$ (чим
		  вони відрізняються від власної індуктивності) і~покажіть на прикладі двох
		  обмоток на спільному осерді, що $L_{12}=L_{21}$. За яких умов ця рівність може
		  порушуватися і~яку роль тут відіграє феромагнітне осердя?
		\item Сформулюйте теорему взаємності. Як теорема спрощує обчислення потоку, як
		  у~прикладі з~двома коаксіальними кільцями?
		\item Узагальніть формулу енергії одного контуру на систему $N$ зв'язаних
		  контурів. Чому саме теорема взаємності дозволяє писати подвійну суму
		  $W=\frac1{2c^2}\sum L_{ik}I_iI_k$, не турбуючись про порядок індексів? Звідки
		  береться нерівність $L_{12}^2\le L_1L_2$ і~що вона означає фізично?
		\item Поясніть принцип дії трансформатора. Що таке коефіцієнт трансформації, чим
		  відрізняються підвищувальний і~знижувальний трансформатори і~чому передавати
		  електроенергію вигідно при високій напрузі?
		\item Виведіть формули для пондеромоторної сили при $I=\const$ та при
		  $\Phi=\const$. Яким режимам електростатики вони відповідають і~чому знаки в~них різні?
		\item Коли енергія $W$ не збігається з~вільною енергією, що визначає роботу поля за $T=\const$? Знайдіть силу,
		  з~якою парамагнітний чи діамагнітний стержень втягується (виштовхується) із~соленоїда.
		\item Знайдіть підйомну силу електромагніта (рис.~\ref{pic:ElectromagnetForce}).
		  Чому вона є силою притягання?
	\end{quizlist}
\end{Questions}

\Closesolutionfile{answer}
